\section{Tuning Changes Head Comparisons, but a Linear Head Stays Ahead}
\label{sec:tuning}

\begin{table}[t]
\centering
\caption{Test accuracy (\%) of three heads on the same CNN, each with its head learning rate
selected by mean validation accuracy over the tested grid (backbone lr is coupled at $0.5\times$
the head lr); mean $\pm$ SD over 3 seeds; bold: highest mean in the row. MLP-8: one hidden
layer of width 8; VQC: the default 8-qubit, 4-layer ring circuit. Differences are in
percentage points, paired over seeds (same training subset and split). The last column
repeats VQC $-$ MLP-8 with both heads at one shared head lr ($10^{-3}$). The linear head's
selected lr lies at the top of its grid in five rows, so its lead may be understated.}
\label{tab:matched}
\small
\setlength{\tabcolsep}{4pt}
\begin{tabular}{llcccccc}
\toprule
& & \multicolumn{3}{c}{Selected lr} & \multicolumn{2}{c}{VQC $-$ (selected lr)} & VQC $-$ MLP-8 \\
\cmidrule(lr){3-5}\cmidrule(lr){6-7}\cmidrule(lr){8-8}
Data & $N_\mathrm{train}$ & Linear & MLP-8 & VQC & Linear & MLP-8 & (lr $10^{-3}$) \\
\midrule
\multirow{3}{*}{MNIST} & 250 & $\mathbf{90.3 \pm 1.9}$ & $78.7 \pm 6.5$ & $76.6 \pm 5.3$ & $-13.7 \pm 5.9$ & $-2.1 \pm 9.7$ & $+10.8 \pm 13.9$ \\
 & 500 & $\mathbf{94.3 \pm 0.2}$ & $86.7 \pm 3.0$ & $90.0 \pm 0.6$ & $-4.3 \pm 0.5$ & $+3.3 \pm 2.6$ & $+6.1 \pm 11.6$ \\
 & 1000 & $\mathbf{96.4 \pm 0.5}$ & $93.9 \pm 0.8$ & $94.6 \pm 0.4$ & $-1.7 \pm 0.9$ & $+0.8 \pm 1.0$ & $-0.5 \pm 3.8$ \\
\midrule
\multirow{3}{*}{CIFAR-10} & 1000 & $\mathbf{59.0 \pm 0.9}$ & $52.3 \pm 1.8$ & $43.1 \pm 3.8$ & $-15.8 \pm 3.8$ & $-9.2 \pm 4.5$ & $-13.6 \pm 1.8$ \\
 & 2000 & $\mathbf{67.6 \pm 0.7}$ & $62.0 \pm 1.0$ & $53.1 \pm 0.9$ & $-14.5 \pm 1.3$ & $-8.9 \pm 1.7$ & $-13.5 \pm 2.0$ \\
 & 4000 & $\mathbf{74.6 \pm 0.5}$ & $70.8 \pm 0.5$ & $67.7 \pm 1.1$ & $-7.0 \pm 1.2$ & $-3.1 \pm 0.9$ & $-10.1 \pm 1.6$ \\
\bottomrule
\end{tabular}
\end{table}


Before turning to representations, we check how much a head comparison depends on the training
recipe. With one shared recipe (head learning rate $10^{-3}$, 25 MNIST epochs; \cref{app:sweep}), the
VQC beat MLP-8 on MNIST by up to 18 points and lost to it on CIFAR-10 by 8--14. Under the protocol of
\cref{sec:setup}, these differences change in size and sign (\cref{tab:matched}): on MNIST, VQC
$-$ MLP-8 becomes $-2.1\pm9.7$, $+3.3\pm2.6$ and $+0.8\pm1.0$ at 250, 500 and 1,000 images, partly
through the protocol itself and partly through per-head tuning (last column); on CIFAR-10 the deficit
shrinks to 3.1--9.2 points. A linear head has the highest mean in all six settings and beats the VQC
for every seed. Every selected rate is interior to its grid.

This audit fixes the comparisons used in the rest of the paper: every head is evaluated at
its own validated learning rate. It also leaves a question. At 4,000 CIFAR-10 images the
VQC trails the linear head by $7.0\pm1.2$ test points (5 seeds: $6.7\pm1.0$). Part of that
gap could come from the VQC as a readout and part from the backbone it trains; a common probe
measures the gap that remains under one readout.
